% ------------------------------------------------------------------------------
% Capítulo 1 · Introducción
% Fuentes: propuesta aprobada del TFG (objetivo, descripción, alcance y EDT),
% docs/tfg/{PLAN,REQUISITOS,DECISIONES}.md.
% ------------------------------------------------------------------------------

Este capítulo presenta la motivación del trabajo, sus objetivos y su alcance,
tal como se fijaron en la propuesta aprobada, y describe la organización del
resto del documento.

\section{Motivación}
\label{sec:intro-motivacion}

Las pequeñas y medianas empresas constituyen la inmensa mayoría del tejido
empresarial europeo~\cite{ec-sme-definition}. Para muchas de ellas, disponer
de una aplicación web propia ---una tienda, una zona de clientes, un sistema de
reservas o una herramienta interna--- ha dejado de ser opcional. El modelo de
\textit{Software as a Service} (SaaS)~\cite{nist-cloud}, en el que la aplicación
se ofrece a través de la red y se paga por uso o por suscripción, es la forma
habitual de construir y explotar este tipo de productos.

Sin embargo, cualquier aplicación SaaS, por sencilla que sea su lógica de
negocio, necesita un conjunto de piezas que se repiten de un proyecto a otro:
registro e inicio de sesión con verificación de correo y segundo factor,
cuentas de usuario y de organización, invitaciones y roles, cobro de
suscripciones, una consola de administración y alguna forma de consultar y
corregir los datos sin escribir consultas a mano. Ninguna de ellas aporta
valor diferencial al producto, pero todas tienen que estar bien hechas: un
error en el control de acceso o en el aislamiento entre clientes compromete
los datos de todas las empresas que usan la aplicación.

Para un equipo pequeño, o para un desarrollador que da servicio a varias
pymes, construir estas piezas desde cero en cada proyecto es lento y
arriesgado. Lo razonable es partir de una base común, ya probada, y dedicar
el esfuerzo a lo que distingue a cada producto. Esa es la idea que motiva este
trabajo: disponer de una \textbf{arquitectura software reutilizable} que
reduzca el tiempo de puesta en marcha de nuevos productos SaaS, facilite la
reutilización de componentes comunes y estandarice buenas prácticas en
seguridad, gestión del acceso, persistencia de datos y administración del
sistema.

\section{Objetivos}
\label{sec:intro-objetivos}

\subsection{Objetivo general}

El objetivo general, tal como figura en la propuesta aprobada, es
\emph{diseñar, implementar y validar una arquitectura software reutilizable
que sirva como base para el desarrollo ágil de aplicaciones SaaS orientadas a
pequeñas y medianas empresas, integrando los componentes esenciales de una
solución moderna: \textit{frontend}, \textit{backend}, sistema de
autenticación, gestión de usuarios y roles, integración de pagos y un módulo
CMS para la administración de contenidos y datos.}

\subsection{Objetivos específicos}

A partir de la descripción y del alcance de la propuesta, el objetivo general
se descompone en los objetivos específicos de la
tabla~\ref{tab:intro-objetivos}. Cada uno se relaciona con los requisitos del
capítulo de Requisitos del Sistema y se retoma en las conclusiones
(capítulo~\ref{cap:conclusiones}) para valorar su grado de cumplimiento.

\begin{table}[htbp]
  \centering
  \small
  \caption{Objetivos específicos del trabajo y requisitos relacionados.}
  \label{tab:intro-objetivos}
  \begin{tabularx}{\textwidth}{@{}lXl@{}}
    \toprule
    \textbf{ID} & \textbf{Objetivo específico} & \textbf{Requisitos} \\
    \midrule
    OE-1 & Analizar el contexto de las aplicaciones SaaS para pymes, identificar los componentes comunes reutilizables y definir los requisitos funcionales y no funcionales de la plataforma base. & Todos \\
    OE-2 & Diseñar una arquitectura modular que separe \textit{frontend}, \textit{backend}, base de datos y servicios auxiliares, con paquetes desacoplados. & RF-01, RF-02, RNF-04 \\
    OE-3 & Diseñar e implementar la autenticación, la autorización y la gestión de usuarios, equipos y roles, aplicando el control de acceso en la propia base de datos. & RF-03--RF-06, RNF-02, RNF-03 \\
    OE-4 & Integrar una pasarela de pagos para la gestión de suscripciones. & RF-07 \\
    OE-5 & Integrar un módulo CMS para la administración de contenidos y de los datos almacenados, con permisos propios y auditoría. & RF-08--RF-11 \\
    OE-6 & Diseñar un modelo de datos multicliente (\textit{multi-tenant}) con relaciones, restricciones y mecanismos de acceso orientados a la reutilización. & RNF-02, RNF-03 \\
    OE-7 & Preparar la plataforma para su reutilización: configuración centralizada, estructura de proyecto reproducible e interfaz en español extensible a otros idiomas. & RF-13, RNF-01, RNF-05 \\
    OE-8 & Validar la solución con pruebas de los principales casos de uso y de la integración entre módulos, y evaluar su reutilización en un escenario de ejemplo. & RNF-01, RNF-06 \\
    OE-9 & Definir el despliegue y elaborar la documentación: memoria, documentación técnica y manuales de instalación y de uso. & RNF-07 \\
    \bottomrule
  \end{tabularx}
\end{table}

\section{Alcance}
\label{sec:intro-alcance}

La propuesta organiza el trabajo en una Estructura de Descomposición del
Trabajo (EDT) con seis paquetes: (1)~análisis del problema y definición de
requisitos; (2)~diseño de la arquitectura software; (3)~diseño del modelo de
datos; (4)~implementación del prototipo funcional; (5)~validación de la
solución; y (6)~documentación. Los entregables comprometidos son la memoria,
el software desarrollado, los ficheros de configuración y despliegue, el
manual de instalación y puesta en marcha y un manual básico de uso. La
relación entre los paquetes de la EDT y las fases reales del proyecto se
detalla en el plan de gestión del proyecto (capítulo~\ref{cap:plan}).

El prototipo debe incluir, como mínimo, gestión de usuarios, autenticación,
persistencia de datos, administración desde un panel CMS e integración con un
flujo de pago. Dentro de ese marco, a lo largo del proyecto se fijaron estas
decisiones de alcance:

\begin{itemize}
  \item \textbf{Plataforma genérica.} PymeKit es una base reutilizable, no
        una aplicación de un sector concreto. Su configuración por defecto
        corresponde a la web de \emph{una} pyme: página pública y blog,
        acceso de clientes y personal, pagos y una consola de administración
        (ADR-005, ADR-020). Los equipos y las suscripciones se conservan
        como módulos que cada instalación activa si los necesita.
  \item \textbf{Un único CMS, de datos.} La propuesta pide un CMS para la
        administración de contenidos y de datos. Se cubre con un CMS que
        trabaja directamente sobre las tablas de la base de datos. El
        contenido editorial (el blog) se guarda también en la base de datos y
        se gestiona desde ese mismo CMS, en lugar de mantener un segundo gestor
        de contenidos basado en ficheros (ADR-004, ADR-017).
  \item \textbf{CMS integrado en la consola.} El CMS no es una aplicación
        aparte: forma parte de la consola de administración de la web, con un
        solo servicio y un solo inicio de sesión (ADR-011).
  \item \textbf{Interfaz solo en español}, sin textos escritos directamente
        en el código, de modo que añadir otro idioma sea un punto de
        extensión documentado (ADR-021).
  \item \textbf{Fuera de alcance:} aplicaciones móviles nativas, la lógica
        de negocio de un sector concreto, la facturación fiscal (emisión de
        facturas con validez legal) y la explotación en producción con
        clientes reales. Los pagos funcionan en el modo de pruebas de la
        pasarela.
\end{itemize}

\paragraph{Aportaciones del trabajo.}
Las aportaciones del trabajo son:

\begin{itemize}
  \item la \textbf{integración del CMS} en la consola de administración de la
        web (ADR-011): montaje de su API dentro del servidor web, pantallas
        implementadas como rutas de la aplicación, sesión única y control de
        acceso coordinado entre los permisos de la plataforma y los del CMS
        (ADR-014, ADR-016) y una sola base de datos para ambos (ADR-012);
  \item la \textbf{unificación de la pila} de librerías de formularios,
        textos y componentes de interfaz en toda la web (ADR-013);
  \item la \textbf{interfaz en español} sobre una infraestructura de
        internacionalización extensible (ADR-021);
  \item la \textbf{consola por áreas de negocio}, con nombres en español
        para las tablas, ayudas contextuales en los formularios y un resumen
        en cada pantalla (ADR-020);
  \item la \textbf{auditoría adversarial y corrección documentada de fallos
        de seguridad}, fijados con pruebas de regresión (ADR-015 y bitácora
        de incidencias);
  \item el \textbf{proceso de desarrollo asistido por IA}, con un arnés de
        directrices, comprobaciones automáticas y revisiones;
  \item la \textbf{configuración reutilizable} centralizada, el
        \textbf{despliegue} y la \textbf{documentación} técnica y de usuario.
\end{itemize}

\section{Glosario de términos}
\label{sec:intro-glosario}

\begin{description}[style=nextline,leftmargin=1.2cm]
  \item[AAL (\textit{Authenticator Assurance Level})] Nivel de garantía de
        una sesión: \texttt{aal1} tras la contraseña o el enlace mágico;
        \texttt{aal2} tras verificar además un segundo factor.
  \item[ADR (\textit{Architecture Decision Record})] Registro breve de una
        decisión de arquitectura con su contexto, alternativas y
        consecuencias~\cite{nygard2011adr}.
  \item[BaaS (\textit{Backend as a Service})] Servicio que ofrece ya
        construidos la base de datos, la autenticación, el almacenamiento y
        una API, de modo que la aplicación no necesita un \textit{backend}
        propio para esas funciones.
  \item[CMS (\textit{Content Management System})] Sistema para gestionar
        contenidos. En este trabajo, un CMS \emph{de datos}: una interfaz para
        consultar y editar las tablas de la base de datos con permisos.
  \item[E2E (\textit{end-to-end})] Prueba que ejercita la aplicación completa
        a través de un navegador, como lo haría un usuario.
  \item[EDT] Estructura de Descomposición del Trabajo (\textit{Work Breakdown
        Structure}).
  \item[Harness (arnés)] Conjunto de directrices, procedimientos y
        comprobaciones automáticas que acota y verifica el trabajo de los
        agentes de IA en el repositorio.
  \item[i18n] Internacionalización: preparar el software para mostrar sus
        textos en varios idiomas sin modificar el código.
  \item[MFA (\textit{Multi-Factor Authentication})] Autenticación con más de
        un factor; aquí, contraseña más un código temporal (TOTP).
  \item[Monorepo] Repositorio único que contiene varias aplicaciones y
        paquetes con dependencias explícitas entre ellos.
  \item[Multi-tenant (multicliente)] Arquitectura en la que una sola
        instalación da servicio a varios clientes (\textit{tenants}) cuyos
        datos deben quedar aislados entre sí.
  \item[pgTAP] Marco de pruebas unitarias para PostgreSQL~\cite{pgtap}.
  \item[Pyme] Pequeña y mediana empresa: menos de 250 personas empleadas y un
        volumen de negocio anual de hasta 50 millones de euros o un balance
        de hasta 43 millones~\cite{ec-sme-definition}.
  \item[RBAC (\textit{Role-Based Access Control})] Control de acceso basado
        en roles a los que se asignan permisos.
  \item[RLS (\textit{Row Level Security})] Mecanismo de PostgreSQL que filtra
        las filas que cada usuario puede leer o modificar mediante
        políticas~\cite{postgresql-rls}.
  \item[SaaS (\textit{Software as a Service})] Modelo en el que el proveedor
        ofrece una aplicación a través de la red sobre una infraestructura
        que gestiona él mismo~\cite{nist-cloud}.
  \item[\textit{Server function}] Función que se ejecuta en el servidor pero
        se invoca desde el cliente como una llamada tipada.
  \item[Super-admin] Administrador de la plataforma, con acceso a la consola
        y a todos los permisos del CMS.
\end{description}

\section{Organización del documento}
\label{sec:intro-organizacion}

La memoria sigue la estructura de la plantilla de la Escuela Superior de
Ingeniería y se divide en cuatro partes:

\begin{itemize}
  \item \textbf{Prolegómeno.} Este capítulo de introducción; el
        capítulo~\ref{cap:contexto}, de antecedentes, con la situación de las
        pymes ante el software SaaS, el estado de la técnica y la
        justificación de la solución; y el capítulo~\ref{cap:plan}, con el
        plan de gestión del proyecto: metodología, tecnologías,
        planificación, organización, costes, riesgos, gestión de la
        configuración y aseguramiento de la calidad.
  \item \textbf{Desarrollo.} Los capítulos de requisitos, análisis, diseño,
        construcción, pruebas y despliegue del sistema, que recorren las
        etapas del proceso de ingeniería con los diagramas, las decisiones y
        los fragmentos de código más significativos.
  \item \textbf{Epílogo.} Las conclusiones (capítulo~\ref{cap:conclusiones}),
        con el grado de cumplimiento de los objetivos, las lecciones
        aprendidas y el trabajo futuro; la información sobre la licencia y la
        bibliografía.
  \item \textbf{Anexos.} El manual del desarrollador
        (anexo~\ref{cap:manual-desarrollador}), el manual de usuario
        (anexo~\ref{cap:manual-usuario}) y la evaluación de la reutilización
        de la plataforma (anexo~\ref{cap:evaluacion}).
\end{itemize}

El software se entrega como un repositorio git con el monorepo completo:
código de la aplicación y del CMS, esquemas, migraciones y pruebas de la base
de datos, pruebas de extremo a extremo, configuración de integración continua
y de despliegue, y la documentación interna del proyecto (plan, requisitos,
decisiones, trazabilidad y bitácora). El repositorio es privado (véase la información sobre
licencia).
